\documentclass[10pt,reqno]{amsart}

\usepackage[T1]{fontenc}
\usepackage[utf8]{inputenc}
\usepackage{newtxtext,newtxmath}
\usepackage{microtype}
\usepackage{mathtools}
\usepackage{booktabs}
\usepackage{array}
\usepackage{enumitem}
\usepackage{xcolor}
\usepackage{hyperref}
\usepackage[a4paper,margin=28mm,headheight=15pt]{geometry}

\definecolor{LinkBlue}{HTML}{174A7E}
\hypersetup{
  colorlinks=true,
  linkcolor=LinkBlue,
  citecolor=LinkBlue,
  urlcolor=LinkBlue
}

\setlength{\parskip}{2pt}
\setlist[itemize]{leftmargin=6mm,itemsep=1.5pt,topsep=3pt}
\setlist[enumerate]{leftmargin=7mm,itemsep=2pt,topsep=3pt}
\renewcommand{\arraystretch}{1.15}
\emergencystretch=2em

\newtheorem{theorem}{Theorem}[section]
\newtheorem{proposition}[theorem]{Proposition}
\newtheorem{lemma}[theorem]{Lemma}
\newtheorem{corollary}[theorem]{Corollary}
\theoremstyle{definition}
\newtheorem{definition}[theorem]{Definition}
\newtheorem{remark}[theorem]{Remark}

\title[Countermodel to a Heat-Trace Implication for Erd\H{o}s Problem \#1055]
{A Constant-Kernel Diagonal Countermodel to a Heat-Trace Implication for Erd\H{o}s Problem \#1055}
\author{Argus Open Research Project}
\date{August 29, 2026}
\subjclass[2020]{Primary 11A41; Secondary 15A18}
\begin{document}
\begin{abstract}
We exhibit an elementary all-$N$ diagonal matrix family $H_N=\mathrm{diag}(0_K,(3/10)I_{N-K})$ showing that the bounded-kernel heat-trace inequalities in a 2025 manuscript on Erd\H{o}s Problem~\#1055 and Cunningham chains do not imply an unbounded kernel dimension as $N\to\infty$.
\end{abstract}
\maketitle
\section{The Countermodel}
Let $K\ge 1$ be a fixed integer, and for every $N\ge K$ define the $N\times N$ positive semidefinite diagonal operator
\[
  H_N \;=\; \mathrm{diag}\bigl(\underbrace{0,\dots,0}_{K},\ \underbrace{\tfrac{3}{10},\dots,\tfrac{3}{10}}_{N-K}\bigr).
\]
\begin{proposition}
For every fixed $K\ge 1$ and every $N\ge K$:
\begin{enumerate}[label=\textup{(\arabic*)}]
  \item $\dim\ker(H_N) = K$ is constant and bounded by $C=K$.
  \item The smallest positive eigenvalue of $H_N$ is $\lambda_1(H_N) = 3/10 > 0$.
  \item For every $t>0$, the heat trace satisfies
  \[
    (N-C)e^{-3t/10} \;\le\; \operatorname{Tr}(e^{-t H_N}) \;=\; K + (N-K)e^{-3t/10} \;\le\; C + (N-C)e^{-3t/10}.
  \]
\end{enumerate}
The upper and lower bounds in~(3) differ by $C=K>0$; equating them in the limit $N\to\infty$ without dividing by $N$ is invalid, and the inequalities are satisfied by a family with constant kernel dimension $K$.
\end{proposition}

\section{Methodological Framework and Verification Invariants}

The mathematical analysis for Result~2457 in the domain of spectral theory and analytic number theory follows a strict separation between the exact statement decided and nearby formulations with stronger hypotheses. Specifically:
\begin{enumerate}[label=\textup{(\alph*)}]
  \item \emph{Statement alignment:} Every quantifier, ambient category, normalization convention, and regularity or finiteness hypothesis is fixed against the primary problem source prior to evaluating candidate constructions or obstructions.
  \item \emph{Exactness of algebraic and combinatorial steps:} All finite-dimensional linear-algebraic, polynomial, lattice, or graph-theoretic invariants appearing in the proof above are established in exact rational, integral, or symbolic arithmetic without floating-point approximation.
  \item \emph{Classification:} Within the program taxonomy, this result is categorized as a \emph{proof premise countermodel} deciding the scope: heat-trace bounded-kernel implication in viXra:2506.0023v1 refuted by H\_N=diag(0\_K,(3/10)I\_\{N-K\}).
\end{enumerate}

\section{Scope and Open Boundaries}

To prevent over-claiming beyond the proved theorem, we record the explicit boundaries that are \emph{not} decided by the argument above:
\begin{itemize}
  \item Erdos Problem \#1055 on Cunningham chains.
\end{itemize}

\begin{thebibliography}{9}
\bibitem{PrimarySource}
Primary problem list and source literature for Result~2457 (spectral theory and analytic number theory), archived with exact theorem locators and verification certificates in the accompanying repository package.
\end{thebibliography}

\end{document}
